\documentclass[11pt]{article}
\usepackage[margin=1in]{geometry}
\usepackage{amsmath,amssymb,amsthm,mathtools,hyperref}
\newtheorem{theorem}{Theorem}
\newtheorem{lemma}{Lemma}
\newtheorem{corollary}{Corollary}
\newcommand{\Tr}{\operatorname{Tr}}
\newcommand{\id}{\operatorname{id}}
\newcommand{\PPT}{\operatorname{PPT}}
\newcommand{\Sep}{\operatorname{Sep}}
\newcommand{\ran}{\operatorname{ran}}
\newcommand{\C}{\mathbb C}
\title{Finite witnesses for measurement information order on PPT and separable preparations}
\author{Astra P3: working derivation, 29 October 2026}
\date{}
\begin{document}
\maketitle

\paragraph{Status and provenance.}
This is a new derivation in this session, not an independently reviewed or
priority-certified result. The accepted Astra N675 undecidability theorem is
used only in Section~\ref{sec:boundary}; its proof is not reopened here.
The resource-preserving Liouville state and resource-interior argument are
explicitly credited to N972. Inverse-action compression has classical and
Petz-quadratic antecedents in N677, N1038 and N676. The classical cutoff was
rediscovered during opportunity selection and is not claimed as new.
The extension developed here is the resource-preserving compression for
arbitrary complex POVM pairs, followed by a computability boundary theorem.
The earlier terminal-calibration construction and its costs remain in
\texttt{../p3\_precision}; they are not assumptions of this paper.

\section{Question and main result}
Let $P=(E_i)_{i=1}^m$ and $Q=(F_j)_{j=1}^n$ be POVMs on
$A=\C^d$. Delete zero effects; all other effects are positive semidefinite,
and each family sums to $I_A$. No reality or simultaneous-diagonalization
assumption is made. The same letters denote the measurement channels.
All logarithms are natural. For a reference $R$, put
$\rho_i=\Tr_A[(E_i\otimes I)\rho]$ and similarly for $Q$.
Let $\mathcal C$ be either the cone of separable positive operators across
$A:R$, or the PPT cone $\{B\geq0:B^{T_A}\geq0\}$.

For fixed finite $c\geq0$, consider:
\begin{align}
 D(P\rho\Vert P\sigma)&\leq cD(Q\rho\Vert Q\sigma),
       &&\rho_R=\sigma_R,                                      \tag{E}\label{eq:E}\\
 D(P\rho\Vert P\sigma)-D(\rho_R\Vert\sigma_R)
   &\leq c\{D(Q\rho\Vert Q\sigma)-D(\rho_R\Vert\sigma_R)\},     \tag{C}\label{eq:C}\\
 g^S_{P\sigma}(PH)&\leq c g^S_{Q\sigma}(QH),
       &&\Tr_A H=0.                                           \tag{S}\label{eq:S}
\end{align}
Here (E) and (C) quantify over resource-class state pairs and finite displayed
entropies; (S) over faithful resource baselines and all Hermitian zero-reference
tangents. The SLD metric is $g^S_a(h)=\langle h,\Omega_a^{-1}h\rangle$,
$\Omega_a=(L_a+R_a)/2$, summed over measurement outcomes.
All finite reference dimensions are initially allowed.

\begin{theorem}[Resource-preserving finite witnesses]\label{thm:main}
The three uniform comparisons (E), (C), (S) are equivalent. If any fails,
there are faithful states $\rho,\sigma\in\mathcal C$ witnessing failure of
(E) with
\[
                   \dim R\leq m+1.
\]
They can be chosen in the interior of the resource class and with
$\rho_A=\sigma_A$ as well as $\rho_R=\sigma_R$.
There is also a violating faithful SLD test in reference dimension at most
$m+1$. Consequently, for rational-complex POVMs and rational $c\geq0$,
these resource-restricted uniform comparisons are decidable exactly.
More precisely, their failure problems belong to $\exists\mathbb R$ under
polynomial-time reductions, so the order problems lie in
$\operatorname{co}\exists\mathbb R\subseteq\mathrm{PSPACE}$.
\end{theorem}

The dimension bound depends on the outcome count of the numerator channel
$P$. It does not bound witness precision, provide a polynomial-time optimizer,
or concern arbitrary quantum-output channels. At $c=1$, (C) is simply the
unrestricted-marginal entropy comparison, since the two reference terms cancel.
PPT is the positive-partial-transpose condition introduced by Peres
\cite{peres}; separability is a strictly stronger restriction in general.

\section{A rectangular inverse-quadratic interface}
Let $B>0$ on $A\otimes R$, with $B\in\mathcal C$; its trace need not be one.
Let $Z:A\longrightarrow A\otimes R$ be an arbitrary complex linear map.
In physical blocks write $Z=(z_{ab})$, $z_{ab}\in R$, and impose
\[
                         \sum_a z_{aa}=0.                       \label{eq:ztrace}\tag{1}
\]
Set
\[
 B_i=\Tr_A[(E_i\otimes I)B],\qquad
 z_i=\sum_{a,b}(E_i)_{ba}z_{ab},\qquad
 q_P(B,Z)=\sum_i z_i^*B_i^{-1}z_i.                              \tag{2}
\]
Every $B_i$ is positive definite: if $B\geq bI$ and $E_i\neq0$, then
$B_i\geq b\Tr(E_i)I>0$. Define $q_Q$ analogously. The rectangular comparison is
\[
                 q_P(B,Z)\leq c q_Q(B,Z)                       \tag{R}\label{eq:R}
\]
for every finite reference, every such $B$, and every $Z$ satisfying (1).

\subsection{The Liouville state}
Use row vectorization, $|X\rangle\!\rangle=\sum_{u,v}X_{uv}|u\rangle|v\rangle$.
For $0\leq s\leq1$, define on $A\otimes R\otimes R'$ ($\dim R'=r$)
\[
 T_s(B)=s(B\otimes I_{R'})+
 (1-s)\operatorname{Swap}_{R,R'}(B^{T_R}\otimes I_{R'})
                 \operatorname{Swap}_{R,R'}.                  \tag{3}
\]
Its trace is $r\Tr B$; normalization is deliberately omitted.
For $B$ PPT, both summands are positive, and their $A$ partial transposes
are positive. For $B$ separable, both summands are separable. Thus $T_s$
preserves either resource cone. Its measurement output is exactly
\[
 (T_s(B))_i=s(B_i\otimes I)+(1-s)(I\otimes B_i^T)
                         =sL_{B_i}+(1-s)R_{B_i}.               \tag{4}
\]
If needed, first replace $B$ by $B+\delta I$. This supplies a strictly
positive identity component throughout (3). Formula (3), particularly its
$s=1/2$ instance, is the resource-preserving construction of N972.

For Hermitian $H$ let $Z_H$ have blocks $|H_{ab}\rangle\!\rangle$.
Then $\Tr_AH=0$ implies (1), and (4) gives
\[
                  q_P(T_{1/2}(B),Z_H)=g^S_{PB}(PH).            \tag{5}
\]
Consequently (R) implies (S).

\subsection{Direct passage to conditional entropy}
This also supplies a direct proof of (R)$\Rightarrow$(C).
For general $Z$ let $b=\Tr_A B$, $z=\sum_a z_{aa}$ and $v=b^{-1}z$.
Put $Z_0=Z-B(I_A\otimes v)$. It satisfies (1). Since every POVM sums to $I$,
\[
 q_P(B,Z_0)=q_P(B,Z)-z^*b^{-1}z,                              \tag{6}
\]
and the same identity holds for $Q$. This is an exact rectangular gauge,
with no assertion that $Z_0$ is a physical finite displacement.

For a faithful resource baseline $\sigma$ and any Hermitian $H$, apply (6)
to $B=T_s(\sigma)$ and $Z=Z_H$. The matrix logarithm derivative satisfies
\[
 g^{\mathrm{BKM}}_a(h)
   =\int_0^1\langle h,[sL_a+(1-s)R_a]^{-1}h\rangle\,ds.       \tag{7}
\]
This follows by diagonalizing $a$ and using
$\int_0^1[su+(1-s)v]^{-1}ds=(\log u-\log v)/(u-v)$.
The same formula applies to the reference block $\sigma_R$.
Integrating the rectangular inequality in $s$ therefore proves the
conditional BKM comparison. If $H=\rho-\sigma$ and
$\sigma_u=\sigma+uH$, the standard Hessian identity
\[
 D(\rho\Vert\sigma)=\int_0^1(1-u)
                 g^{\mathrm{BKM}}_{\sigma_u}(H)\,du           \tag{8}
\]
proves (C). Resource-class interiors are convex. Boundary pairs with all
displayed entropies finite follow by common white-noise regularization and
joint convexity/lower semicontinuity. Thus (R)$\Rightarrow$(C)$\Rightarrow$(E).

\subsection{Vacuum realization}
Conversely fix $B,Z$ with (1). Normalize $B$ to trace one; this multiplies
both quadratic forms by the same positive scalar and preserves a violation.
White-noise regularization preserves any strict violation for sufficiently
small positive noise. Hence assume $B$ is interior PPT or has a decomposition
$B=B_{\rm sep}+\kappa_0 I$, $B_{\rm sep}$ separable, $\kappa_0>0$.
On $\widetilde R=R\oplus\C|v\rangle$ set
\[
 \Gamma_\epsilon=
   \begin{pmatrix}(1-\epsilon)B&0\\0&\epsilon I_A/d\end{pmatrix},
 \qquad
 X=\begin{pmatrix}0&Z\\Z^*&0\end{pmatrix}.                    \tag{9}
\]
Block notation in (9) orders the reference sectors; the physical system
$A$ is unchanged. Both partial traces of $X$ vanish. The output at $E_i$ is
\[
 \Gamma_i=\operatorname{diag}((1-\epsilon)B_i,\epsilon a_i),
 \quad X_i=\begin{pmatrix}0&z_i\\z_i^*&0\end{pmatrix},
 \quad a_i=\Tr E_i/d>0.                                      \tag{10}
\]
The baseline is interior in the chosen resource class. In the separable case
it has a global identity component of any size
$0<\kappa<\min\{(1-\epsilon)\kappa_0,\epsilon/d\}$, with a separable remainder.
Thus $\Gamma_\epsilon+tX$ remains faithful and in the same class for all
sufficiently small real $t$.

Diagonalizing $B_i$ gives
\begin{align}
 \lim_{\epsilon\downarrow0}g^S_{\Gamma_i}(X_i)&=4z_i^*B_i^{-1}z_i,
                                                                  \tag{11}\\
 \lim_{\epsilon\downarrow0}
   \frac{g^{\mathrm{BKM}}_{\Gamma_i}(X_i)}{2\log(1/\epsilon)}
                       &=z_i^*B_i^{-1}z_i.                       \tag{12}
\end{align}
For fixed $\epsilon$, $D(\Gamma_i+tX_i\Vert\Gamma_i)
=t^2 g^{\mathrm{BKM}}_{\Gamma_i}(X_i)/2+O(t^3)$.
Applying (E), taking $t\to0$, then $\epsilon\downarrow0$, and finally removing
the optional white noise proves (R). Applying (S) to (9), then taking
$\epsilon\downarrow0$, also proves (R). This establishes the equivalence of
(E), (C), (S), (R), independently of the compression below.

\section{One-sided compression preserves the resource}
\begin{lemma}[Inverse-action compression]\label{lem:compression}
If $q_P(B,Z)>c q_Q(B,Z)$, $c\geq0$, there is another such test with
reference dimension at most $m$, and with $B$ in the same resource cone.
\end{lemma}
\begin{proof}
Let $u_i=B_i^{-1}z_i$, and let $V:\C^k\to R$ be an isometry onto
$\operatorname{span}\{u_1,\ldots,u_m\}$, where $1\leq k\leq m$.
Define
\[
 B'=(I_A\otimes V^*)B(I_A\otimes V),\qquad
 Z'=(I_A\otimes V^*)Z.                                      \tag{13}
\]
Local reference compression preserves separability and PPT, and $B'>0$.
The zero-reference constraint survives. Write $u_i=V\widehat u_i$. Then
$B'_i\widehat u_i=V^*B_i u_i=V^*z_i=z'_i$, whence
\[
                    q_P(B',Z')=q_P(B,Z).                     \tag{14}
\]
For any $D>0$ and vector $w$, the variational identity
\[
 w^*D^{-1}w=\max_v\{2\Re(w^*v)-v^*Dv\}
\]
shows that compressing $D,w$ to any subspace can only decrease this value.
Applied separately to the $Q$ outcomes, it gives
$q_Q(B',Z')\leq q_Q(B,Z)$. Since $c\geq0$, the strict violation survives.
\end{proof}

Apply the vacuum realization (9)--(12) to the compressed test after optional
arbitrarily small identity regularization. For small finite $\epsilon$ the
BKM inequality is strictly violated; for a subsequent small finite $t$ the
actual entropy pair violates (E). Its reference dimension is $k+1\leq m+1$.
Both input marginals agree because both partial traces of $X$ vanish.
Equation (11) gives the same dimension bound for SLD. This proves the
finite-witness part of Theorem~\ref{thm:main}. This is an existence and
description construction; compression is not asserted to be a deterministic
physical operation preserving normalized statistics of an unknown state.

\section{An exact finite decision procedure and its costs}
It suffices to decide whether (R) fails with $k=m$. A smaller witness can
be embedded and padded with a zero-tangent reference block. Normalize
$\Tr B=1$ and $\|Z\|_F=1$. Use real and imaginary parts as real variables,
and introduce vectors $u_i,v_j$ satisfying
\[
 B_i u_i=z_i,\qquad B_j^Q v_j=z_j^Q,\qquad
 \sum_i\Re(z_i^*u_i)-c\sum_j\Re((z_j^Q)^*v_j)>0.              \tag{15}
\]
These are polynomial constraints. It is enough to require $B>0$ and
$B^{T_A}>0$ in the PPT case: any strict violation on the resource boundary
survives identity regularization. Positive definiteness is a finite system
of positive leading principal minors. For the separable case, conic
Caratheodory in the real Hermitian space gives a decomposition into at most
$(dm)^2$ positive rank-one product terms. Writing their local vectors as
variables makes separability a finite existential polynomial condition.
One can require an additional positive identity component without changing
the answer, again by continuity of a strict violation.

Real quantifier elimination decides (15), including all equality cases, and
semialgebraic sampling produces real-algebraic witness data on the negative
side. Rational approximation after strict interior regularization is also
possible. The PPT representation uses on the order of
$(dm)^2+d^2m+(m+n)m$ real scalar degrees of freedom before harmless constant
factors; a product decomposition adds up to order $(dm)^2(d+m)$ variables.
This is an exact finite algorithm, not a polynomial-runtime assertion.
Quantifier elimination, algebraic-number descriptions and conditioning can
be very expensive. Dense emitted state matrices have $d^2(m+1)^2$ entries,
each with its own precision cost. A finite cutoff is not a fixed-size
numerical grid and does not certify conclusions from finitely many random
tests.

\paragraph{Polynomial-size existential description.}
The preceding encoding can be made polynomial in the explicit rational input
length. Avoid expanding determinants: introduce complex lower-triangular
matrices $L,M$ with strictly positive real diagonal and impose
$B=LL^*$, $B^{T_A}=MM^*$, entrywise. These are quadratic equations and strict
linear inequalities. For separability instead impose
$B=\delta I+\sum_{\ell=1}^{(dm)^2}
u_\ell u_\ell^*\otimes v_\ell v_\ell^*$, $\delta>0$; the equations have degree
at most four. Normalization, (1), (15) and the auxiliary inverse equations
also have polynomial size. Thus failure is a polynomial reduction to the
existential theory of the reals. Its PSPACE upper bound is classical
\cite{canny}; it does not imply polynomial time or short rational witnesses.
No matching hardness classification is established here.

\begin{corollary}[Algebraic optimal coefficients]
For rational-complex POVMs, the least constant $c\geq0$ in any of (E), (C),
(S) on either resource class is either $+\infty$ or a real algebraic number.
An exact description by polynomial equations and an isolating interval is
computable. Finite attainment by an entropy pair is not asserted.
\end{corollary}
\begin{proof}
Keep $c$ as one free real variable in the finite polynomial description.
Quantifier elimination makes its failure set semialgebraic over the rationals.
The admissible set is an intersection of closed half-lines, since each
individual condition is $q_P-cq_Q\leq0$ with $q_Q\geq0$. It is therefore
empty or $[c_*,\infty)$ for $c_*\geq0$. A finite endpoint of a rational
semialgebraic subset of the line is real algebraic. Elimination and univariate
root isolation provide the stated exact description. Its degree and output
bit length need not be small.
\end{proof}

\subsection{Explicit finite precision for the vacuum step}
Let a normalized, interior top test have strict gap
$g=q_P-cq_Q>0$. Let $a>0$ be a lower bound on every output eigenvalue of
$B$, and let $q_0>0$ lower-bound every $\Tr E_i/d$ and $\Tr F_j/d$.
All these output eigenvalues and probabilities are at most one. Set
\[
 L=\log\frac{2}{\min(a,q_0)},\qquad
 S=\sum_i\|z_i\|^2+c\sum_j\|z_j^Q\|^2.
\]
For $\epsilon\leq a/4$, $\gamma=\log(1/\epsilon)$, direct scalar algebra gives
\[
 \left|\frac{1}{\gamma}
 \frac{\log((1-\epsilon)\lambda)-\log(\epsilon q)}
 {(1-\epsilon)\lambda-\epsilon q}-\frac1\lambda\right|
 \leq\frac{2L}{\gamma a}+\frac{4\epsilon}{a^2}.               \tag{16}
\]
Thus it is sufficient to choose a rational $\epsilon=2^{-k}$ satisfying
\[
 \gamma\geq\frac{8LS}{ag},\qquad
 \epsilon\leq\min\{a/4,ga^2/(16S)\}.                         \tag{17}
\]
The normalized BKM gap is then at least $g/2$. This explicitly exposes the
cost of small spectral floors and small strict gaps. No uniform polynomial
bit-precision bound is inferred.

For a fixed faithful output block $A$ with tangent $Y$, put $\nu=\lambda_{\min}A$.
When $A+sY\geq A/2$ along the segment, differentiating the resolvent Hessian
yields
\[
 \left|D(A+tY\Vert A)-\frac{t^2}{2}g_A^{\mathrm{BKM}}(Y)\right|
 \leq\frac{2|t|^3}{3}\frac{\|Y\|_{\rm op}\|Y\|_F^2}{\nu^2}, \tag{18}
\]
for the trace-zero blocks in (10). Let $b>0$ be the BKM gap and $M$ the sum
of the right-side cubic coefficients, with weight $c$ on $Q$.
Taking $0<t\leq b/(4M)$, together with the positivity bounds below, yields
an actual entropy gap at least $t^2b/4$.
For PPT it suffices additionally to take
\[
 t\leq\tfrac12\min\left\{
 \frac{\lambda_{\min}\Gamma}{\|X\|_{\rm op}},
 \frac{\lambda_{\min}\Gamma^{T_A}}{\|X^{T_A}\|_{\rm op}}
 \right\}.                                                   \tag{19}
\]
Zero denominators in (19) impose no restriction.

For separability there is an elementary explicit alternative. If
$\Gamma-\kappa I$ is separable on dimensions $d,D$, then
\[
                   |t|dD\|X\|_F\leq\kappa/2                 \tag{20}
\]
suffices. Indeed expand $X$ in orthonormal local Hermitian bases whose
operator norms are at most one. The coefficient $\ell_1$ norm is at most
$dD\|X\|_F$. For Hermitian contractions $U,V$, both
$I\otimes I\pm U\otimes V$ are separable, as follows from
$\tfrac12[(I+U)\otimes(I\pm V)+(I-U)\otimes(I\mp V)]$.
This gives a finite product decomposition and a global positive margin.
Sharper separable balls are prior work \cite{gurvits}; optimality is not needed.

Unknown measurement devices must first be acquired; the theorem takes explicit
effects as input. Exact equality is not resolved by finite-precision tomography.
For a fixed test, if an output $B_i$ changes by $\Delta B_i$ and $z_i$ by
$\Delta z_i$, with $\|\Delta B_i\|\leq a_i/2$ and $B_i\geq a_iI$, then
\[
 |\widetilde z_i^*\widetilde B_i^{-1}\widetilde z_i-z_i^*B_i^{-1}z_i|
 \leq \frac{2\|z_i\|^2\|\Delta B_i\|}{a_i^2}
       +\frac{4\|z_i\|\|\Delta z_i\|+2\|\Delta z_i\|^2}{a_i}.\tag{21}
\]
Summing (21) gives a sufficient error budget for an already acquired strict
witness. It gives no uniform noise tolerance at the exact order boundary.

\section{The computability boundary}\label{sec:boundary}
This section explicitly assumes the programme-accepted N675 theorem:
for variable-dimensional rational-complex strictly positive POVM pairs,
unrestricted complete SLD/relative-entropy order is coRE-complete; failures
are RE; and a computable $V=0$ versus $V\geq1$ reduction exists for the
bounded unrestricted entropy defect. Reference marginals need not agree.
No re-proof or stronger interpretation of that accepted result is used.

Let $\mathsf U$ be the set of such pairs obeying unrestricted entropy order
$D(P\rho\Vert P\sigma)\leq D(Q\rho\Vert Q\sigma)$ for all finite references.
Let $\mathsf T$ be the set obeying the same order for all PPT pairs.
Theorem~\ref{thm:main} makes $\mathsf T$ decidable, and
$\mathsf U\subseteq\mathsf T$.

\begin{corollary}[Undecidability remains inside the PPT-valid promise]
The restriction of $\mathsf U$ to the decidable promise $\mathsf T$ is
coRE-complete and has no computable sufficient reference cutoff for complete
SLD testing.
In particular, there exist rational strictly positive POVM pairs for which
all PPT preparations obey the order but some NPT preparation violates it.
\end{corollary}
\begin{proof}
If $\mathsf T\setminus\mathsf U$ were empty, $\mathsf U$ would be decidable.
Choose a fixed $x_0$ in this difference. This is an effective existence:
enumerate rational POVM pairs, decide membership in $\mathsf T$, and dovetail
the RE search for a full SLD violation. A pair $x_0$ is eventually found.
Now map an arbitrary pair $x$ to $x$ when $x\in\mathsf T$, and to $x_0$
otherwise. This total computable map lands in $\mathsf T$ and preserves
membership in $\mathsf U$. A computable SLD cutoff on this decidable promise,
combined with finite-dimensional SLD quantifier elimination, would decide it.
\end{proof}

This is not an explicitly acquired numerical example of $x_0$. Its search has
not been run, and neither its size nor search time is claimed bounded feasibly.
The consequence concerns exact computation from descriptions, not the ability
to prepare or discover the NPT witness physically.
No separate decidability assumption is made for fixed-reference entropy
optimization; the no-cutoff deduction is explicitly an SLD deduction.

\begin{lemma}[PPT order tensorizes for measurement pairs]
If $(P_1,Q_1)$ and $(P_2,Q_2)$ belong to $\mathsf T$, then so does
$(P_1\otimes P_2,Q_1\otimes Q_2)$.
\end{lemma}
\begin{proof}
Start with a state PPT across $A_1A_2:R$. After measuring $A_2$ with an effect
$G$, the conditional state $\tau$ on $A_1R$ satisfies
\[
 \tau^{T_{A_1}}
   =\Tr_{A_2}[(G^T\otimes I)\rho^{T_{A_1A_2}}]\geq0.
\]
Keeping the classical outcome as part of the reference preserves PPT.
Thus one may compare $P_1$ with $Q_1$ after $P_2$, and then compare $P_2$
with $Q_2$ after $Q_1$. The usual hybrid entropy argument proves the claim.
The same argument preserves separability. General complete-order
tensorization is prior \cite{hirche}; this proof checks the restricted cone.
\end{proof}

The fixed bad pair $x_0$ can therefore be tensor-amplified, using an actually
found positive entropy gap, to a fixed PPT-valid pair with defect at least one.
Apply the preceding decidable filter to the accepted N675 gap reduction.
It follows that its fixed additive approximation obstruction also persists
under the PPT-valid promise. Tensor product descriptions may grow enormously;
this reduction is computable, not claimed efficient.

\begin{theorem}[Any positive flagged neighborhood restores the full order]
Fix any rational $0<\eta<1$. Restrict each pair of hypotheses to states lying
within trace distance $\eta$ of the same product state on $A:R$, allowing
arbitrary finite references. The resulting uniform unrestricted-marginal
entropy order is exactly $\mathsf U$. It is therefore coRE-complete, even
under the PPT-valid promise. The corresponding supremum defect has the
same fixed additive approximation obstruction after computable tensor
amplification. Here trace distance is $\tfrac12\|\cdot\|_1$.
\end{theorem}
\begin{proof}
Given any violating faithful pair $\rho,\sigma$ on $A\otimes R$, choose a
faithful signal state $\tau$, a faithful reference state $\omega$, and
$0<w<\eta$. On the flagged reference $\C|v\rangle\oplus R$ put
\[
 \widetilde\rho=(1-w)\tau\otimes|v\rangle\langle v|\oplus w\rho,
 \quad
 \widetilde\sigma=(1-w)\tau\otimes|v\rangle\langle v|\oplus w\sigma.
                                                               \tag{22}
\]
Both are faithful and within distance at most $w$ of the same product state
$\tau\otimes[(1-w)|v\rangle\langle v|\oplus w\omega]$.
The orthogonal flag and identical weights give exactly
\[
 D(P\widetilde\rho\Vert P\widetilde\sigma)
 -D(Q\widetilde\rho\Vert Q\widetilde\sigma)
 =w\{D(P\rho\Vert P\sigma)-D(Q\rho\Vert Q\sigma)\}.          \tag{23}
\]
Thus every strict full-order failure occurs in this neighborhood. The reverse
implication is immediate. For a fixed additive gap, start from the accepted
$V=0$ versus $V\geq1$ instances on the PPT-valid promise, take $k$ tensor
copies with $wk\geq2$, and then flag as above. The zero case remains zero;
the other case has restricted supremum at least two. PPT order is preserved
by the tensorization lemma. Every dimension and emitted tensor entry belongs
to the construction cost.
\end{proof}

This is a sharp change at an exact resource boundary, not a claim that generic
depolarization preserves hard witnesses. The construction uses an accessible
orthogonal flag and a rare, possibly very complicated branch. Without prior
tensor amplification its entropy gap shrinks by exactly $w$. The theorem
does not assert hardness for equal-reference hypotheses in the unrestricted
NPT class; the accepted N675 entropy interface allows unequal marginals.
One can also make negativity arbitrarily small: flag dilution multiplies
$(\|\rho^{T_A}\|_1-1)/2$ by $w$. This yields the same exact-order statement
for any positive negativity allowance, but no fixed additive-gap claim is
made under that different quantitative constraint.

\section{Attacks, failures and scope}
The initial conjecture that classical-device complete order might inherit
N675 undecidability was abandoned: N677 and N1038 already provide finite
classical certificates/cutoffs. A duplicated-classical-label construction
separates numerator and denominator coefficients only on a chosen slice;
it does not encode the arbitrary CP baseline cone of N675. This failed route
helped expose the actual resource-preserving compression above.

The same argument does not apply to unrestricted NPT baselines: (3) can lose
positivity. Nor does it apply automatically to quantum-output channels:
the desired Liouville action would require a second copy of the transmitted
output, which a local application of a general channel does not provide.
Full rank is never used as a substitute for interior separability.
The compression preserves a strict comparison, not normalized experimental
payoff, sampling probability, or the cost of preparing the input.

The independent comparison-learning opening is saved in
\texttt{APN\_DIFFERENTIAL\_SAMPLING.tex}. Its new operation does not constitute
an efficient next-coordinate learner. The low-defect many-body route was
screened but not used as a theorem premise. No new upload, archive access,
publication, external review, or physical realization is asserted.

\begin{thebibliography}{9}
\bibitem{peres} A. Peres, Separability Criterion for Density Matrices,
Phys. Rev. Lett. 77, 1413 (1996),
\href{https://doi.org/10.1103/PhysRevLett.77.1413}{doi:10.1103/PhysRevLett.77.1413}.
\bibitem{gurvits} L. Gurvits and H. Barnum, Largest separable balls around the
maximally mixed bipartite quantum state, Phys. Rev. A 66, 062311 (2002),
\href{https://arxiv.org/abs/quant-ph/0204159}{arXiv:quant-ph/0204159}.
\bibitem{hirche} C. Hirche, C. Rouze and D. Stilck Franca,
On contraction coefficients, partial orders and approximation of capacities
for quantum channels, Quantum 6, 862 (2022),
\href{https://doi.org/10.22331/q-2022-11-28-862}{doi:10.22331/q-2022-11-28-862}.
\bibitem{canny} J. F. Canny, Some Algebraic and Geometric Computations in
PSPACE, STOC 1988, pp. 460--467,
\href{https://doi.org/10.1145/62212.62257}{doi:10.1145/62212.62257};
\href{https://www2.eecs.berkeley.edu/Pubs/TechRpts/1988/6041.html}{author's institutional report}.
\bibitem{astra} Astra source revision
\texttt{1d03a0f1942512a1132478741373b73511a30935}:
N972 governing \texttt{povm\_resource\_cone\_order\_equivalence.md},
source-byte SHA256
\texttt{7398907c550937eb2c128d2e0d6970b71f49e163d319d8f97d868b005accb533};
exact locators and hashes are recorded in \texttt{SOURCE\_AND\_CLAIM\_NOTES.json}.
N677 and N1038 supply the classical antecedents; N675 supplies only the
programme-accepted undecidability premises stated in Section~\ref{sec:boundary}.
\end{thebibliography}
\end{document}
